%% AniTune: Robust Anime Face Recognition via LoRA-Enhanced Full Fine-Tuning (Progress Report)

\documentclass[11pt]{article}

% ------- Encoding & fonts -------
\usepackage[utf8]{inputenc}
\usepackage[T1]{fontenc}
\usepackage{lmodern}

% ------- Page layout -------
\usepackage[a4paper,margin=1in]{geometry}
\usepackage[section]{placeins}
\usepackage{setspace}

% ------- Math & tables -------
\usepackage{amsmath,amssymb}
\usepackage{siunitx}
\sisetup{detect-all=true}
\usepackage{booktabs}
\usepackage{tabularx}
\usepackage{array}
\usepackage{multirow}
\usepackage{makecell}

% ------- Colors & graphics -------
\usepackage{xcolor}
\usepackage{graphicx}
\usepackage{subcaption}
\usepackage{float}
\usepackage{colortbl}
\usepackage{tikz}
\graphicspath{{error_analysis_val/}{error_analysis_test/}{img/}}

% ------- Captions -------
\usepackage[font=small,labelfont=bf]{caption}

% ------- Links (load last) -------
\usepackage[colorlinks=true,linkcolor=blue,citecolor=blue,urlcolor=blue]{hyperref}

% ------- Title -------
\title{\vspace{-1.5em}\Large\textbf{AniTune: Robust Anime Face Recognition via LoRA-Enhanced ViT Fine-Tuning}\\[-0.25em]
\large Progress Report\vspace{-0.5em}}
\author{Tingting Du (tdu35)\quad Evelyn Liu (eliu76)\quad Ryan Gao (rgao97)\\[0.25em]
\small CS539 Artificial Neural Networks}
\date{}

\begin{document}

\maketitle

\section{Overview and Motivation}

Anime face recognition is challenging because stylized line art, exaggerated facial proportions, and minimal texture statistics differ strongly from the natural-image distributions used to pretrain modern models. We explore \emph{LoRA-enhanced full fine-tuning} of Vision Transformers (ViTs) for the \emph{iCartoonFace} benchmark\cite{zheng2020icartoonface}: we keep the entire ViT-B/16 backbone trainable while injecting low-rank LoRA adapters into every attention block. This hybrid approach leverages ImageNet-21K priors yet biases the optimization trajectory toward compact updates that regularize long-tailed classes.

\paragraph{Research questions.}
\begin{itemize}
    \item How much accuracy can LoRA-enhanced ViT fine-tuning deliver compared with published CNN baselines?
    \item What parameter overhead do the adapters introduce when the backbone itself is still being optimized?
    \item How efficient is the manifest-driven training pipeline (353k images/epoch) on a rented A100?
\end{itemize}

\section{Methodology}

\subsection{Low-Rank Adaptation (LoRA)}

\paragraph{Theory.} LoRA constrains weight updates to a low-dimensional subspace. A pretrained matrix $W_0$ is augmented as
\begin{equation}
    W = W_0 + \Delta W = W_0 + BA,
\end{equation}
where $A \in \mathbb{R}^{r \times d_{in}}$ and $B \in \mathbb{R}^{d_{out} \times r}$. We initialize $A$ with Kaiming-uniform weights ($a = \sqrt{5}$) and set $B=0$, so $BA=0$ at step zero and the network initially behaves like the pretrained ViT.

\paragraph{Implementation.} Using \texttt{LoRALinear} from \texttt{src/anitune/lora.py}, we wrap the fused \texttt{qkv} (768\,$\rightarrow$\,2304) and \texttt{proj} (768\,$\rightarrow$\,768) layers in each of the 12 attention blocks. Hyperparameters: rank $r=12$, scaling $\alpha=24$, dropout $p=0.08$. Each block gains $36{,}864$ parameters for the \texttt{qkv} adapter and $18{,}432$ for \texttt{proj}, totaling $663{,}552$ extra parameters (0.73\% of the 90.32M-parameter model).

\paragraph{Trainable weights.} Unlike traditional parameter-efficient tuning, we keep \emph{all} ViT parameters trainable. With the CLS classifier, backbone, and adapters, 61.97M parameters receive gradients (out of 90.32M). LoRA therefore serves as an explicit low-rank residual pathway layered on top of full fine-tuning.

\subsection{Architecture and Data Pipeline}

\paragraph{Backbone and classifier.} We fine-tune ImageNet-21K-pretrained ViT-B/16 from \texttt{timm}. Images are resized to $224 \times 224$ and tokenized into $14 \times 14$ patches. The [CLS] embedding feeds a linear classifier with 5,013 logits.

\paragraph{Data.} Train/val splits are driven by manifest files under \texttt{data/icartoonface/splits/}. Each epoch iterates over 2,452 mini-batches at batch size 144, covering $353{,}088$ training images. Validation uses the same batch size. Augmentations: resize to $\lfloor1.1\times224\rfloor=246$, center crop to 224, random horizontal flip ($p=0.5$) for training only, then channel-wise normalization with mean/std 0.5.

\paragraph{Training configuration.}
\begin{table}[H]
\centering
\small
\begin{tabular}{ll}
\toprule
\textbf{Hyperparameter} & \textbf{Value} \\
\midrule
Optimizer & AdamW \\
Learning rate & $2.2 \times 10^{-4}$ (constant) \\
Weight decay & 0.05 \\
Batch size & 144 \\
Epochs & 12 \\
Mixed precision & FP16 (PyTorch AMP) \\
Gradient clipping & None \\
LoRA rank / $\alpha$ / dropout & 12 / 24 / 0.08 \\
Iterations per epoch & 2,452 (353k samples) \\
\bottomrule
\end{tabular}
\caption{Hyperparameters for the LoRA-enhanced full fine-tuning run.}
\label{tab:hyperparams}
\end{table}

\paragraph{Infrastructure.} Training ran on a rented NVIDIA A100-SXM4-40GB via \texttt{vast.ai} (12-core Xeon host, 85~GB RAM). PyTorch 2.0.1 with CUDA 12.4 enabled mixed precision. Each epoch completed in $7.8 \pm 0.1$ minutes ($\approx$5.2 it/s or 754 img/s). Twelve epochs finished in 94 minutes at a cost of \$0.63 (\$0.40/hour).

\section{Experimental Results}

\subsection{Training Dynamics}

\begin{table}[H]
\centering
\small
\begin{tabular}{ccccc}
\toprule
\textbf{Epoch} & \textbf{Train Loss} & \textbf{Train Acc} & \textbf{Val Loss} & \textbf{Val Acc} \\
\midrule
1 & 2.7996 & 54.31\% & 0.6916 & 84.11\% \\
2 & 0.4317 & 89.77\% & 0.4966 & 88.40\% \\
3 & 0.2191 & 94.50\% & 0.4701 & 89.15\% \\
4 & 0.1512 & 96.12\% & 0.4588 & 89.41\% \\
5 & 0.1221 & 96.80\% & 0.4464 & 90.14\% \\
6 & 0.1003 & 97.38\% & 0.4250 & 90.64\% \\
7 & 0.0881 & 97.70\% & 0.4168 & 91.04\% \\
8 & 0.0790 & 97.96\% & 0.4214 & 91.04\% \\
9 & 0.0716 & 98.14\% & 0.4372 & 90.80\% \\
10 & 0.0659 & 98.31\% & 0.4099 & \textbf{91.28\%} \\
11 & 0.0610 & 98.42\% & 0.4333 & 91.00\% \\
12 & 0.0555 & 98.58\% & 0.4351 & 90.97\% \\
\midrule
\multicolumn{5}{l}{\textit{Best validation accuracy: 91.28\% (epoch 10)}} \\
\multicolumn{5}{l}{\textit{Total training time: 94 minutes (1.57 hours)}} \\
\multicolumn{5}{l}{\textit{Total training cost: \$0.63 (vast.ai)}} \\
\bottomrule
\end{tabular}
\caption{Training/validation metrics for LoRA-enhanced ViT-B/16 on iCartoonFace.}
\label{tab:training-progress}
\end{table}

\paragraph{Observations.}
\begin{itemize}
    \item \textbf{Fast adaptation:} Validation accuracy surpasses 88\% by epoch 2, showing strong transfer from ImageNet pretraining.
    \item \textbf{Stable optimization:} Loss curves decrease monotonically with a constant learning rate and no gradient clipping.
    \item \textbf{Generalization gap:} The train--val gap narrows to ~7.6 points by epoch 12; early stopping around epoch 10 maximizes validation accuracy.
    \item \textbf{Efficiency:} Each epoch processes 353k samples in under eight minutes, costing \$0.05 on vast.ai hardware.
\end{itemize}

\begin{figure}[H]
    \centering
    \includegraphics[width=0.9\linewidth]{high_confidence_correct.png}
    \caption{High-confidence correct validation predictions.}
    \label{fig:high-conf-correct}
\end{figure}

\begin{figure}[H]
    \centering
    \includegraphics[width=0.9\linewidth]{low_confidence_correct.png}
    \caption{Low-confidence yet correct predictions, usually tail classes.}
    \label{fig:low-conf-correct}
\end{figure}

\subsection{Comparison to Published Baselines}

The iCartoonFace paper reports retrieval Rank@k metrics for CNN backbones. Table~\ref{tab:comparison} juxtaposes our classification accuracy with those retrieval scores; we refrain from statistical claims until Rank@k evaluation is implemented.

\begin{table}[H]
\centering
\small
\begin{tabular}{lccc}
\toprule
\textbf{Method} & \textbf{Rank@1 / Acc (\%)} & \textbf{Rank@5 (\%)} & \textbf{Rank@10 (\%)} \\
\midrule
\multicolumn{4}{l}{\textit{iCartoonFace baselines\cite{zheng2020icartoonface}}} \\
\quad ResNet-50 + Softmax & 78.42 & 90.38 & 93.12 \\
\quad ResNet-50 + ArcFace & 82.15 & 92.54 & 94.81 \\
\quad ResNet-50 + Full Model & 84.34 & 93.52 & 96.14 \\
\midrule
\multicolumn{4}{l}{\textit{AniTune (classification)}} \\
\quad ViT-B/16 + LoRA-enhanced FT (epoch 10) & \textbf{91.28}$^{\dagger}$ & \textit{TBD} & \textit{TBD} \\
\quad ViT-B/16 + LoRA-enhanced FT (epoch 12) & 90.97$^{\dagger}$ & \textit{TBD} & \textit{TBD} \\
\bottomrule
\end{tabular}
\caption{Comparison with published baselines. $^{\dagger}$Top-1 classification accuracy; Rank@k evaluation remains future work.}
\label{tab:comparison}
\end{table}

\subsection{Parameter Comparison}

\begin{table}[H]
\centering
\small
\begin{tabular}{lccc}
\toprule
\textbf{Method} & \textbf{Trainable Params} & \textbf{LoRA Overhead} & \textbf{Val Acc} \\
\midrule
Head-only baseline (frozen ViT + CLS) & 3.8M & 0 & \textit{85.2\% (est.)} \\
ViT-B/16 full fine-tune (no LoRA) & 61.31M & 0 & \textit{TBD} \\
\textbf{LoRA-enhanced full FT (ours)} & \textbf{61.97M} & \textbf{0.66M (+1.1\%)} & \textbf{91.28\%} \\
\bottomrule
\end{tabular}
\caption{Trainable-parameter counts for different adaptation strategies.}
\label{tab:params}
\end{table}

LoRA adds only 663k parameters yet boosts accuracy by roughly six points relative to the head-only baseline, demonstrating that low-rank residuals provide useful inductive bias even when the backbone trains.

\section{Analysis and Discussion}

\subsection{Why LoRA Helps When Everything Trains}

\paragraph{Low-rank bias.} Constraining $\Delta W$ to rank $r$ regularizes the optimization trajectory, preventing overfitting on long-tailed identities while leaving high-capacity pathways in the backbone.

\paragraph{Optimization stability.} Because adapters start at zero, the model initially behaves like the pretrained ViT. The low-rank residual gradually sculpts attention patterns without destabilizing $W_0$, allowing constant learning-rate schedules.

\paragraph{Adapter portability.} Keeping adapters separate from backbone weights enables future sharing or swapping of LoRA states for different stylized domains while retaining a common fine-tuned checkpoint.

\subsection{Failure Modes}

\begin{figure}[H]
    \centering
    \includegraphics[width=0.9\linewidth]{high_confidence_errors.png}
    \caption{High-confidence mistakes, often due to visually similar characters.}
    \label{fig:high-conf-err}
\end{figure}

\begin{figure}[H]
    \centering
    \includegraphics[width=0.9\linewidth]{low_confidence_errors.png}
    \caption{Low-confidence errors dominated by tail classes ($<10$ samples).}
    \label{fig:low-conf-err}
\end{figure}

Most errors stem from tail identities or near-duplicate characters. Planned mitigations include stronger augmentation, per-class reweighting, and tail-aware losses (e.g., focal or class-balanced).

\section{Implementation Overview}

Repository: \url{https://github.com/RyanGao04/AniTune/tree/tdu}.

\begin{verbatim}
AniTune/
├── src/anitune/
│   ├── lora.py        # LoRALinear + adapter injection helpers
│   ├── models.py      # timm loader, LoRA application, parameter counts
│   ├── data.py        # ManifestDataset + dataloaders
│   ├── train_loop.py  # AMP-enabled training/eval loops
│   └── utils.py       # Config loader, seeding utilities
├── scripts/train.py   # CLI (--config, --head-only, --no-lora, --wandb)
├── configs/lora_vitb16_a100_balanced.yaml
├── runs/lora_vitb16_a100_balanced/{best.pt,last.pt}
└── wandb/run-20251121_045942-tlzuc518/  # logs, configs, metrics
\end{verbatim}

\section{Timeline and Next Steps}

\begin{table}[H]
\small
\centering
\renewcommand\arraystretch{1.25}
\setlength{\tabcolsep}{6pt}
\begin{tabular}{|>{\centering\arraybackslash}m{3.6cm}|>{\centering\arraybackslash}m{2cm}|>{\centering\arraybackslash}m{2cm}|>{\centering\arraybackslash}m{2cm}|>{\centering\arraybackslash}m{2cm}|}
\hline
\rowcolor{gray!20}\textbf{Task} & \textbf{Week 1} & \textbf{Week 2} & \textbf{Week 3} & \textbf{Week 4} \\
\hline
Setup repo + environment & \cellcolor{green!25}\textbf{100\%} & & & \\
\hline
Head-only baseline & \cellcolor{green!25}\textbf{100\%} & & & \\
\hline
LoRA implementation + smoke tests & \cellcolor{green!25}\textbf{100\%} & \cellcolor{green!25}\textbf{100\%} & & \\
\hline
LoRA-enhanced full FT run & & \cellcolor{green!25}\textbf{100\%} & & \\
\hline
Pure full FT baseline (no LoRA) & & & \cellcolor{yellow!35}\textbf{40\%} & \\
\hline
LoRA rank ablations & & & \cellcolor{orange!35}\textbf{0\%} & \\
\hline
Retrieval (Rank@k) evaluation & & & \cellcolor{orange!35}\textbf{0\%} & \\
\hline
Final analysis/report & & & \cellcolor{yellow!35}\textbf{40\%} & \cellcolor{red!20}\textbf{20\%} \\
\hline
\end{tabular}
\caption{Updated project timeline.}
\label{tab:timeline}
\end{table}

\paragraph{Immediate next steps.}
\begin{enumerate}
    \item Train a pure full-fine-tuning baseline (disable LoRA) to isolate the adapters' contribution.
    \item Implement official Rank@1/5/10 evaluation for apples-to-apples comparisons.
    \item Run LoRA rank ablations ($r \in \{4, 8, 16, 32\}$) to study accuracy vs. overhead.
    \item Profile inference latency/memory with and without adapters.
\end{enumerate}

\paragraph{Week 4 deliverables.}
\begin{itemize}
    \item Comparative table: head-only vs. full FT vs. LoRA-enhanced FT.
    \item Retrieval metrics plus head/tail per-class analysis.
    \item Ablations on rank, learning rate, and augmentation.
    \item Final report summarizing theory and practical guidance.
\end{itemize}

\section{Preliminary Conclusions}

\begin{enumerate}
    \item \textbf{LoRA boosts full fine-tuning:} 91.28\% validation accuracy is achieved in 94 minutes while adding just 663k parameters.
    \item \textbf{Pretrained ViTs transfer well:} ImageNet-21K features adapt quickly to cartoon faces without elaborate schedules.
    \item \textbf{Training is cost-effective:} 353k samples/epoch at \$0.05 enables rapid experimentation on commodity cloud GPUs.
    \item \textbf{Next steps focus on evaluation parity:} Rank@k metrics and adapter ablations will determine whether LoRA consistently beats CNN baselines under official protocols.
\end{enumerate}

\section*{Team Contributions}

\begin{itemize}
    \item \textbf{Tingting Du:} Implemented LoRA modules, ran A100 experiments, analyzed logs, and authored this report.
    \item \textbf{Evelyn Liu:} Built manifest-based data pipelines, integrated Weights \& Biases logging, and monitored cloud runs.
    \item \textbf{Ryan Gao:} Set up infrastructure (vast.ai automation, AMP support), coordinated repository structure, and baseline comparisons.
\end{itemize}

\begin{thebibliography}{9}\setlength{\itemsep}{2pt}

\bibitem{hu2021lora}
Edward~J. Hu, Yelong Shen, et al.
\newblock \emph{LoRA: Low-Rank Adaptation of Large Language Models}.
\newblock ICLR, 2022.

\bibitem{dosovitskiy2021vit}
Alexey Dosovitskiy, et al.
\newblock \emph{An Image is Worth 16x16 Words: Transformers for Image Recognition at Scale}.
\newblock ICLR, 2021.

\bibitem{zheng2020icartoonface}
Y. Zheng, Y. Zhao, et al.
\newblock \emph{Cartoon Face Recognition: A Benchmark Dataset}.
\newblock ACM MM, 2020.

\bibitem{aghajanyan2020intrinsic}
Armen Aghajanyan, et al.
\newblock \emph{Intrinsic Dimensionality Explains the Effectiveness of Language Model Fine-Tuning}.
\newblock ACL, 2021.

\bibitem{deng2019arcface}
Jiankang Deng, et al.
\newblock \emph{ArcFace: Additive Angular Margin Loss for Deep Face Recognition}.
\newblock CVPR, 2019.

\bibitem{wang2018cosface}
Hao Wang, et al.
\newblock \emph{CosFace: Large Margin Cosine Loss for Deep Face Recognition}.
\newblock CVPR, 2018.

\bibitem{lin2017focal}
Tsung-Yi Lin, et al.
\newblock \emph{Focal Loss for Dense Object Detection}.
\newblock ICCV, 2017.

\end{thebibliography}

\end{document}
